\documentclass[10pt,a4paper]{amsart}
\usepackage{amsmath,amssymb,amsthm,graphicx}
% Retained minimal preamble for the Stacks Project chapter flat.tex (More on Flatness).
% Verbatim excerpt of the chapter preamble: Xy-pic setup (native preamble lines 14-18) and the native
% theorem declarations with the mathematical macros (native preamble lines 154-217). Copied mechanically;
% nothing added, removed or redefined. The wrapper class and the amsmath/amssymb/amsthm/graphicx base
% packages are supplied by the assembler, so the class/externaldocument part of the native preamble is omitted.
% For commutative diagrams we use Xy-pic
\usepackage[all]{xy}

% We use 2cell for 2-commutative diagrams.
\xyoption{2cell}
% Theorem environments.
%
\theoremstyle{plain}
\newtheorem{theorem}[subsection]{Theorem}
\newtheorem{proposition}[subsection]{Proposition}
\newtheorem{lemma}[subsection]{Lemma}

\theoremstyle{definition}
\newtheorem{definition}[subsection]{Definition}
\newtheorem{example}[subsection]{Example}
\newtheorem{exercise}[subsection]{Exercise}
\newtheorem{situation}[subsection]{Situation}

\theoremstyle{remark}
\newtheorem{remark}[subsection]{Remark}
\newtheorem{remarks}[subsection]{Remarks}

\numberwithin{equation}{subsection}

% Macros
%
\def\lim{\mathop{\mathrm{lim}}\nolimits}
\def\colim{\mathop{\mathrm{colim}}\nolimits}
\def\Spec{\mathop{\mathrm{Spec}}}
\def\Hom{\mathop{\mathrm{Hom}}\nolimits}
\def\Ext{\mathop{\mathrm{Ext}}\nolimits}
\def\SheafHom{\mathop{\mathcal{H}\!\mathit{om}}\nolimits}
\def\SheafExt{\mathop{\mathcal{E}\!\mathit{xt}}\nolimits}
\def\Sch{\mathit{Sch}}
\def\Mor{\mathop{\mathrm{Mor}}\nolimits}
\def\Ob{\mathop{\mathrm{Ob}}\nolimits}
\def\Sh{\mathop{\mathit{Sh}}\nolimits}
\def\NL{\mathop{N\!L}\nolimits}
\def\CH{\mathop{\mathrm{CH}}\nolimits}
\def\proetale{{pro\text{-}\acute{e}tale}}
\def\etale{{\acute{e}tale}}
\def\QCoh{\mathit{QCoh}}
\def\Ker{\mathop{\mathrm{Ker}}}
\def\Im{\mathop{\mathrm{Im}}}
\def\Coker{\mathop{\mathrm{Coker}}}
\def\Coim{\mathop{\mathrm{Coim}}}

% Boxtimes
%
\DeclareMathSymbol{\boxtimes}{\mathbin}{AMSa}{"02}

%
% Macros for moduli stacks/spaces
%
\def\QCohstack{\mathcal{QC}\!\mathit{oh}}
\def\Cohstack{\mathcal{C}\!\mathit{oh}}
\def\Spacesstack{\mathcal{S}\!\mathit{paces}}
\def\Quotfunctor{\mathrm{Quot}}
\def\Hilbfunctor{\mathrm{Hilb}}
\def\Curvesstack{\mathcal{C}\!\mathit{urves}}
\def\Polarizedstack{\mathcal{P}\!\mathit{olarized}}
\def\Complexesstack{\mathcal{C}\!\mathit{omplexes}}
% \Pic is the operator that assigns to X its picard group, usage \Pic(X)
% \Picardstack_{X/B} denotes the Picard stack of X over B
% \Picardfunctor_{X/B} denotes the Picard functor of X over B
\def\Pic{\mathop{\mathrm{Pic}}\nolimits}
\def\Picardstack{\mathcal{P}\!\mathit{ic}}
\def\Picardfunctor{\mathrm{Pic}}
\def\Deformationcategory{\mathcal{D}\!\mathit{ef}}

% --- appended 2026-10-12 (Chief): the project own definitions of the reference/slogan/history
% display environments, copied mechanically from the shared Stacks preamble lines 8-12
% (\usepackage{verbatim} plus the three \newenvironment lines, which the chapter preamble keeps
% verbatim from the project preamble). Nothing added, removed or redefined: these environments
% carry no mathematical content and the project itself typesets them as comments. Without these
% three lines a counted statement whose own text uses them fails with Environment slogan undefined.
\usepackage{verbatim}
\newenvironment{reference}{\comment}{\endcomment}
\newenvironment{slogan}{\comment}{\endcomment}
\newenvironment{history}{\comment}{\endcomment}
\usepackage{cleveref}
\crefname{originallemma}{original lemma}{original lemmas}
\Crefname{originallemma}{Original lemma}{Original lemmas}
\makeatletter
\newlabel{algebra-lemma-finite-projective-descends}{{00O1}{source}{}{sourcepointer-1}{}}
\newlabel{algebra-lemma-finite-projective-descends@cref}{{[originallemma][0][]00O1}{[1][0][]source}{}{}{}}
\makeatother
\crefname{originallemma}{original lemma}{original lemmas}
\Crefname{originallemma}{Original lemma}{Original lemmas}
\makeatletter
\newlabel{equation-free-at-generic-point-fibre}{{05MQ}{source}{}{sourcepointer-2}{}}
\newlabel{equation-free-at-generic-point-fibre@cref}{{[originallemma][0][]05MQ}{[1][0][]source}{}{}{}}
\makeatother
\crefname{originallemma}{original lemma}{original lemmas}
\Crefname{originallemma}{Original lemma}{Original lemmas}
\makeatletter
\newlabel{lemma-induction-step}{{05I3}{source}{}{sourcepointer-3}{}}
\newlabel{lemma-induction-step@cref}{{[originallemma][0][]05I3}{[1][0][]source}{}{}{}}
\makeatother
\crefname{originallemma}{original lemma}{original lemmas}
\Crefname{originallemma}{Original lemma}{Original lemmas}
\makeatletter
\newlabel{limits-lemma-descend-module-flat-finite-presentation}{{05LY}{source}{}{sourcepointer-4}{}}
\newlabel{limits-lemma-descend-module-flat-finite-presentation@cref}{{[originallemma][0][]05LY}{[1][0][]source}{}{}{}}
\makeatother
\crefname{originallemma}{original lemma}{original lemmas}
\Crefname{originallemma}{Original lemma}{Original lemmas}
\makeatletter
\newlabel{limits-lemma-descend-opens}{{01Z4}{source}{}{sourcepointer-5}{}}
\newlabel{limits-lemma-descend-opens@cref}{{[originallemma][0][]01Z4}{[1][0][]source}{}{}{}}
\makeatother
\crefname{originallemma}{original lemma}{original lemmas}
\Crefname{originallemma}{Original lemma}{Original lemmas}
\makeatletter
\newlabel{modules-lemma-finite-type-surjective-on-stalk}{{01B8}{source}{}{sourcepointer-6}{}}
\newlabel{modules-lemma-finite-type-surjective-on-stalk@cref}{{[originallemma][0][]01B8}{[1][0][]source}{}{}{}}
\makeatother
\crefname{originallemma}{original lemma}{original lemmas}
\Crefname{originallemma}{Original lemma}{Original lemmas}
\makeatletter
\newlabel{more-morphisms-theorem-openness-flatness}{{0399}{source}{}{sourcepointer-7}{}}
\newlabel{more-morphisms-theorem-openness-flatness@cref}{{[originallemma][0][]0399}{[1][0][]source}{}{}{}}
\makeatother
\crefname{originallemma}{original lemma}{original lemmas}
\Crefname{originallemma}{Original lemma}{Original lemmas}
\makeatletter
\newlabel{situation-free-at-generic-points}{{05MP}{source}{}{sourcepointer-8}{}}
\newlabel{situation-free-at-generic-points@cref}{{[originallemma][0][]05MP}{[1][0][]source}{}{}{}}
\makeatother
\title{The functor of points where a module is free at the generic points of the fibres is an fpqc sheaf and limit preserving}
\author{The Stacks Project Authors}
\date{stacks-project:flat.tex:local-snapshot-20261011; GNU FDL 1.2 or later}
\begin{document}
\maketitle
\noindent\textit{Editorial note.} All delivered mathematical text is the author's own, copied from the pinned Stacks Project chapter file flat.tex; no statement or proof step was rewritten, reordered or removed.\par\medskip
\noindent\textit{Editorial note.} Preamble provenance: the chapter's own retained minimal preamble (sources/stacks-flat/retained-preamble.tex, lines 1-85) is a verbatim excerpt of the native Xy-pic setup and of the native theorem declarations and mathematical macros, followed by the project's own reference/slogan/history display-environment definitions copied mechanically from the shared preamble; the wrapper class and the amsmath/amssymb/amsthm/graphicx base packages are supplied by the assembler, and extra\_packages is left empty because everything needed is already in the retained preamble. No macro was invented or redefined.\par\medskip
\noindent\textit{Editorial note.} Same-chapter pointers are located at their \textbackslash{}label lines inside the pinned chapter file; cross-chapter pointers are located in the project's own tag list staged in the snapshot (sources/stacks-chow/tags.txt), at the verbatim lines carrying their canonical tags.\par\medskip
\noindent\textit{Editorial note.} The author's closing sentence at source line 6267, 'Some details omitted.', is retained verbatim at the end of the counted proof, as is the author's own footnote attached to the fpqc covering in the first line of the proof. The sheaf argument and the limit preserving reduction are both carried out in the span; only the repeated verification of the same reduction is left to the reader.\par\medskip
\noindent\textit{Allowed standard prerequisites (editorial).}\begin{itemize}
\item The situation defining the functor of points where the module is free of a fixed rank at the generic point of the fibre. Same chapter flat.tex, situation and equation labelled situation-free-at-generic-points and equation-free-at-generic-point-fibre (statement, proof line 6180; used at proof lines 6197, 6241 and 6266). Fixes what H\_p is and what condition on the pullback of F the proof has to establish; the two parts of the proof are then organised around that condition.
\item Descent of finite projective modules along a flat local homomorphism, and generation of a finite type module on a neighbourhood. Algebra, label algebra-lemma-finite-projective-descends (proof line 6209); Modules, label modules-lemma-finite-type-surjective-on-stalk (proof line 6213). The first propagates the local freeness of the pullback to the limit point through the flat local map of stalks; the second produces an affine neighbourhood of the generic point and a surjection from a free module onto the restriction of the module.
\item The induction step lemma and the openness of flatness, with the limit descent lemmas. Same chapter flat.tex, label lemma-induction-step (proof line 6225); More on Morphisms, label more-morphisms-theorem-openness-flatness (proof line 6246); Limits, labels limits-lemma-descend-opens (proof line 6251) and limits-lemma-descend-module-flat-finite-presentation (proof line 6255). The induction step converts the isomorphism on every fibre into injectivity after tensoring the global sections with the local rings, which upgrades the displayed surjection to an isomorphism; openness of flatness and the two descent lemmas reduce the limit preserving question to a stage of the inverse system.
\end{itemize}\medskip
\ifdefined\hypertarget\hypertarget{sourcepointer-1}{}\fi
\noindent\textit{Source reference.} Original lemma 00O1: sources/stacks-chow/tags.txt, source lines 879--879. Navigation pointer for \textbackslash{}ref\{algebra-lemma-finite-projective-descends\} used inside the counted span (line 6209); the argument that consumes it is retained in the span.\par\medskip
\ifdefined\hypertarget\hypertarget{sourcepointer-2}{}\fi
\noindent\textit{Source reference.} Original lemma 05MQ: sources/stacks-flat/pinned-original-source.tex, source lines 6160--6160. Navigation pointer for \textbackslash{}ref\{equation-free-at-generic-point-fibre\} used inside the counted span (line 6197); the argument that consumes it is retained in the span.\par\medskip
\ifdefined\hypertarget\hypertarget{sourcepointer-3}{}\fi
\noindent\textit{Source reference.} Original lemma 05I3: sources/stacks-flat/pinned-original-source.tex, source lines 2356--2356. Navigation pointer for \textbackslash{}ref\{lemma-induction-step\} used inside the counted span (line 6225); the argument that consumes it is retained in the span.\par\medskip
\ifdefined\hypertarget\hypertarget{sourcepointer-4}{}\fi
\noindent\textit{Source reference.} Original lemma 05LY: sources/stacks-chow/tags.txt, source lines 7252--7252. Navigation pointer for \textbackslash{}ref\{limits-lemma-descend-module-flat-finite-presentation\} used inside the counted span (line 6255); the argument that consumes it is retained in the span.\par\medskip
\ifdefined\hypertarget\hypertarget{sourcepointer-5}{}\fi
\noindent\textit{Source reference.} Original lemma 01Z4: sources/stacks-chow/tags.txt, source lines 2579--2579. Navigation pointer for \textbackslash{}ref\{limits-lemma-descend-opens\} used inside the counted span (line 6251); the argument that consumes it is retained in the span.\par\medskip
\ifdefined\hypertarget\hypertarget{sourcepointer-6}{}\fi
\noindent\textit{Source reference.} Original lemma 01B8: sources/stacks-chow/tags.txt, source lines 1717--1717. Navigation pointer for \textbackslash{}ref\{modules-lemma-finite-type-surjective-on-stalk\} used inside the counted span (line 6213); the argument that consumes it is retained in the span.\par\medskip
\ifdefined\hypertarget\hypertarget{sourcepointer-7}{}\fi
\noindent\textit{Source reference.} Original lemma 0399: sources/stacks-chow/tags.txt, source lines 4264--4264. Navigation pointer for \textbackslash{}ref\{more-morphisms-theorem-openness-flatness\} used inside the counted span (line 6246); the argument that consumes it is retained in the span.\par\medskip
\ifdefined\hypertarget\hypertarget{sourcepointer-8}{}\fi
\noindent\textit{Source reference.} Original lemma 05MP: sources/stacks-flat/pinned-original-source.tex, source lines 6146--6146. Navigation pointer for \textbackslash{}ref\{situation-free-at-generic-points\} used inside the counted span (line 6180); the argument that consumes it is retained in the span.\par\medskip
\begin{lemma}
\label{lemma-free-at-generic-points}
In Situation \ref{situation-free-at-generic-points}.
\begin{enumerate}
\item The functor $H_p$ satisfies the sheaf property for the fpqc topology.
\item If $\mathcal{F}$ is of finite presentation, then functor $H_p$ is
limit preserving.
\end{enumerate}
\end{lemma}

\begin{proof}
Let $\{T_i \to T\}_{i \in I}$ be an fpqc\footnote{It is quite easy to
show that $H_p$
is a sheaf for the fppf topology using that flat morphisms of finite
presentation are open. This is all we really need later on. But it is kind
of fun to prove directly that it also satisfies the sheaf condition
for the fpqc topology.} covering of schemes over $S$.
Set $X_i = X_{T_i} = X \times_S T_i$ and denote $\mathcal{F}_i$ the
pullback of $\mathcal{F}$ to $X_i$. Assume that $\mathcal{F}_i$ satisfies
(\ref{equation-free-at-generic-point-fibre}) for all $i$.
Pick $t \in T$ and let $\xi_t \in X_T$ denote the generic point of $X_t$.
We have to show that $\mathcal{F}$ is free in a neighbourhood of $\xi_t$.
For some $i \in I$ we can find a $t_i \in T_i$ mapping to $t$.
Let $\xi_i \in X_i$ denote the generic point of $X_{t_i}$, so that
$\xi_i$ maps to $\xi_t$. The fact that $\mathcal{F}_i$ is free of rank
$p$ in a neighbourhood of $\xi_i$ implies that
$(\mathcal{F}_i)_{x_i} \cong \mathcal{O}_{X_i, x_i}^{\oplus p}$
which implies that
$\mathcal{F}_{T, \xi_t} \cong \mathcal{O}_{X_T, \xi_t}^{\oplus p}$
as $\mathcal{O}_{X_T, \xi_t} \to \mathcal{O}_{X_i, x_i}$ is flat, see
for example
Algebra, Lemma \ref{algebra-lemma-finite-projective-descends}.
Thus there exists an affine neighbourhood $U$ of $\xi_t$ in $X_T$
and a surjection
$\mathcal{O}_U^{\oplus p} \to \mathcal{F}_U = \mathcal{F}_T|_U$, see
Modules, Lemma \ref{modules-lemma-finite-type-surjective-on-stalk}.
After shrinking $T$ we may assume that $U \to T$ is surjective. Hence
$U \to T$ is a smooth morphism of affines with geometrically irreducible
fibres. Moreover, for every $t' \in T$ we see that the induced map
$$
\alpha :
\mathcal{O}_{U, \xi_{t'}}^{\oplus p}
\longrightarrow
\mathcal{F}_{U, \xi_{t'}}
$$
is an isomorphism (since by the same argument as before the module on the right
is free of rank $p$). It follows from
Lemma \ref{lemma-induction-step}
that
$$
\Gamma(U, \mathcal{O}_U^{\oplus p})
\otimes_{\Gamma(T, \mathcal{O}_T)} \mathcal{O}_{T, t'}
\longrightarrow
\Gamma(U, \mathcal{F}_U)
\otimes_{\Gamma(T, \mathcal{O}_T)} \mathcal{O}_{T, t'}
$$
is injective for every $t' \in T$. Hence we see the surjection $\alpha$ is an
isomorphism. This finishes the proof of (1).

\medskip\noindent
Assume that $\mathcal{F}$ is of finite presentation.
Let $T = \lim_{i \in I} T_i$ be a directed limit of affine $S$-schemes
and assume that $\mathcal{F}_T$ satisfies
(\ref{equation-free-at-generic-point-fibre}).
Set $X_i = X_{T_i} = X \times_S T_i$ and denote $\mathcal{F}_i$ the
pullback of $\mathcal{F}$ to $X_i$.
Let $U \subset X_T$ denote the open subscheme of points where
$\mathcal{F}_T$ is flat over $T$, see
More on Morphisms, Theorem \ref{more-morphisms-theorem-openness-flatness}.
By assumption every generic point of every fibre is a point of $U$, i.e.,
$U \to T$ is a smooth surjective morphism with geometrically irreducible
fibres. We may shrink $U$ a bit and assume that $U$ is quasi-compact.
Using
Limits, Lemma \ref{limits-lemma-descend-opens}
we can find an $i \in I$ and a quasi-compact open $U_i \subset X_i$
whose inverse image in $X_T$ is $U$. After increasing $i$ we may
assume that $\mathcal{F}_i|_{U_i}$ is flat over $T_i$, see
Limits, Lemma \ref{limits-lemma-descend-module-flat-finite-presentation}.
In particular, $\mathcal{F}_i|_{U_i}$ is finite locally free
hence defines a locally constant
rank function $\rho : U_i \to \{0, 1, 2, \ldots \}$.
Let $(U_i)_p \subset U_i$ denote the open and closed
subset where $\rho$ has value $p$. Let $V_i \subset T_i$ be the
image of $(U_i)_p$; note that $V_i$ is open and quasi-compact.
By assumption the image of $T \to T_i$ is contained in $V_i$.
Hence there exists an $i' \geq i$ such that $T_{i'} \to T_i$ factors
through $V_i$ by
Limits, Lemma \ref{limits-lemma-descend-opens}.
Then $\mathcal{F}_{i'}$ satisfies (\ref{equation-free-at-generic-point-fibre})
as desired. Some details omitted.
\end{proof}

\end{document}
